% Generated draft: requires mathematical and source review.
\documentclass[11pt]{article}
\usepackage[margin=25mm]{geometry}
\usepackage{fontspec}
\setmainfont{Latin Modern Roman}
\usepackage{amsmath,amssymb,mathtools,unicode-math}
\setmathfont{Latin Modern Math}
\usepackage{enumitem,needspace,xurl,hyperref}
\setlength{\parindent}{0pt}
\setlength{\parskip}{5pt}
\setlength{\emergencystretch}{3em}
\newcommand{\N}{\mathbb{N}}
\newcommand{\Z}{\mathbb{Z}}
\newcommand{\Q}{\mathbb{Q}}
\newcommand{\R}{\mathbb{R}}
\newcommand{\C}{\mathbb{C}}
\newcommand{\F}{\mathbb{F}}
\newcommand{\A}{\mathbb{A}}
\newcommand{\PP}{\mathbb{P}}
\newcommand{\E}{\mathbb{E}}
\newcommand{\eps}{\varepsilon}
\newcommand{\dd}{\,\mathrm{d}}
\begin{document}
\section*{UnsolvedMath ID: 10400207}
\textbf{Source problem code:} catalogue entry
\subsection{Definitions and mathematical statement}
Benedetti asks whether the torsions discussed in the surrounding section depend only on the \((pL)\)-homotopy-immersion class of a \((pL)\)-knot. If \(\mathcal K_{pL}(W)\) is the intended space of such knots, \(h\) its projection to the intended homotopy-through-immersions classes, and \(\mathcal T\) one of the source's torsions, the proposed factorization is
\[h(K)=h(K')\quad\Longrightarrow\quad\mathcal T(K)=\mathcal T(K'),\]
with the appropriate identification of torsion targets and auxiliary data. Equivalently \(\mathcal T\) should factor through \(h\). In a homotopy through immersions, intermediate maps need not be embeddings; any \((pL)\) constraints must hold as specified by the missing definition. The symbols \((pL)\), the allowed knot category and the torsions are not defined here. In particular the letters must not silently be expanded as ``piecewise linear'' from the research-summary paraphrase; the Legendrian context requires checking the intended category.

The following background fixes a category \(C\) and a \(C\)-homotopy-immersion class \(\alpha\), and denotes by \(F(\alpha)\) the finite-type invariants of the \(C\)-isotopy classes lying in \(\alpha\). For a class \(\alpha_0\) of Legendrian knots, it sets \(\alpha_1=f_1(\alpha_0)\) and \(\alpha_2=f_2(\alpha_1)\), giving pullbacks
\[F(\alpha_2)\xrightarrow{f_2^*}F(\alpha_1)\xrightarrow{f_1^*}F(\alpha_0).\]
Thus an invariant at stage \(i\) pulls back to one at stage \(i-1\). The record reports that [130], under certain hypotheses on \(W\), for instance the stated integral-homology-sphere setting, makes \(f_1^*\circ f_2^*\) bijective. It contrasts examples in [380] where this composition is not surjective: Legendrian finite-type invariants can distinguish certain knots isotopic as framed knots. It further says that in these examples \(f_2^*\) is already not surjective, so \((pL)\)-finite-type invariants can distinguish certain \((pL)\)-knots isotopic as framed knots. These reports concern all finite-type invariants and their forgetful maps, and do not by themselves decide the factorization question for the particular torsions.

The excerpt ends by referring to a following conjecture compatible with these results, but does not supply that conjecture. It is not reconstructed as an additional assertion here.
\subsection{Short English statement}
Do the specified torsions forget all information beyond a knot's \((pL)\)-homotopy-immersion class? The surrounding finite-type comparison maps and their reported bijective and non-surjective cases are retained, while the undefined knot category and torsions remain explicit gaps.
\subsection{Sources}
\begin{itemize}
\item Local UnsolvedMath registry, permanent ID \(10400207\), exact source code \texttt{catalogue entry}.
\item R. Benedetti, attribution and authorship of Section 12.1, in Ohtsuki, \emph{Problems on invariants of knots and 3-manifolds} (2002), Question 12.1, PDF page 158: \url{https://msp.org/gtm/2002/04/gtm-2002-04-024s.pdf}. Registry access 2026-07-29; pdf-font-structure extraction.
\item Internal references [130] (bijectivity under hypotheses) and [380] (non-surjectivity examples); full entries are absent.
\end{itemize}
\end{document}
